
\subsection{Enhanced Section Strains}\label{sec:enhan-section}
In principle, any deficiency of the standard kinematic ansatz may be addressed by enriching the warping expansion in \eqref{eq:finite-alignment-field}, i.e., increasing $n_w$ and introducing additional warping amplitudes $\bm{\alpha}(x)$. This strategy is systematic, but it is often excessive in practice because each added mode promotes an additional one-dimensional field that must be interpolated along $\mathcal{C}$ and solved globally. 
Moreover, many of the warping features that materially affect the response are not independent kinematic degrees of freedom: they are predominantly local cross-sectional distortions whose amplitude is governed by the existing generalized strains $\bm{e}$ and therefore arises in proportion to them. 
This motivates an alternative viewpoint in which the beam-level balance laws are left unchanged, but the section strain state entering the constitutive model is enriched. The enhanced formulation introduced below follows the classical enhanced strain concept: additional section modes are appended to the compatible strain field and their amplitudes are determined locally by a stationarity/orthogonality condition with respect to the admissible (Saint-Venant) stress space, so that no additional global degrees of freedom are introduced and the construction applies equally to geometrically exact and linearized members of the beam family.

With this motivation, the section strain field $\hat{\StrainVector}$ is augmented as follows.


The section strain field \(\hat{\StrainVector}\) may be enhanced in a manner analogous to the three-field variational formulation of elasticity \citep{simo1986variational,simo1990class,simo1992geometrically}:
\begin{SaveEquation}{eq:enhan-strain}
\bar{\bm{E}} = \bar{\bm{\varepsilon}}\symbun\FrameBasis
\qquad\text{with}\qquad
\bar{\StrainVector}(\reflenx,\bm{r}) =  \hat{\StrainVector}(\bm{e}) + \EnhanShapes(\bm{r}) \bm{\EnhanParam}(x)
\end{SaveEquation}
where \(\hat{\bm{\varepsilon}}\) is the compatible strain vector derived from the displacement field \eqref{eq:finite-alignment-field}, and $\EnhanShapes(\bm{r}) \bm{\EnhanParam}(x)$ is the enhanced part of the strain field. 

Let $\operatorname{A}\left[\mathcal{E}\right] \subset \mathscr{E}$ be the space of compatible linear strain tensors generated by the space $\mathcal{E}$ of rod strain fields \(\bm{e}\); i.e. 
$$
\begin{aligned}
\operatorname{B^*}\left[\mathcal{E}\right]
&\coloneq \left\{
\hat{\StrainVector} \symbun \FrameBasis \in \mathscr{E} 
\bigm| 
\hat{\StrainVector}
=\TotalStrainMatrix(\xi,\hat{\rho}) \bm{e},
\quad \bm{e} \in \mathcal{E}
\right\}
\end{aligned}
$$

Define the finite dimensional space generated by assumed strain interpolations \(\EnhanShapes\bigl(\xi,\hat{\rho}\bigr)\)
\begin{equation*}
\tilde{\mathscr{E}}
\coloneq 
\left\{
\Enhan{\StrainVector} \symbun \FrameBasis \in \tilde{\mathscr{E}} \mid 
\tilde{\bm{\varepsilon}}(\xi,\hat{\rho})=\EnhanShapes(\xi, \hat{\rho}) \bm{\EnhanParam} , \text { and } \bm{\EnhanParam} \in \mathbb{R}^{\EnhanDim}\right\}
\label{eq:SR16}
\end{equation*}
where $\bm{\EnhanParam} \in \mathbb{R}^{n_a}$ are $\EnhanDim$ local parameters. 

The following conditions are considered:
\begin{enumerate}[label=\roman*]
    \item Stability. 
    The enhanced strain space $\tilde{\mathscr{E}}$ and the compatible strain space \(\operatorname{B^*}\left[\mathcal{E}\right]\) have null intersection
    $$
    \tilde{\mathscr{E}} \cap \operatorname{B^*}\left[\mathcal{E}\right]=\{0\}
    $$

\item Consistency. 
    The subspace of admissible stresses, denoted by $\StressSpaceH \subset L$, is $L_2$ -orthogonal to the space $\tilde{\mathscr{E}}$ of enhanced strain fields. 
    This requires \(\residn(\bm{\EnhanParam}) = \mathbf{0}\) with
    \begin{SaveEquation}{eq:def-enhan-residual}
    \residn(\bm{\EnhanParam}) \coloneq \int_{\Section} \EnhanShapes^{\top} \StressVector \left(\operatorname{A}[\bm{e}] + \EnhanShapes \bm{\EnhanParam} \right) \dA
    \end{SaveEquation}
\end{enumerate}
Condition (i) is essential. Condition (ii) is ideal and justified as follows:
\begin{itemize}
    \item When...
\end{itemize}
In traditional assumed strain formulations a third criteria is required
which asserts that the space of admissible nominal stress fields contains piece-wise constant functions, and implies \citep{simo1990class}:
\begin{equation}
\int_{\Section \times \FrameBody} \EnhanShapes \dA \, \dd \xi = \mathbf{0}
\label{eq:condiii}
\end{equation}
However, such a condition would not make sense in the present context. 


\subsection{Linearization}

Linearizing \eqref{eq:stress-resultants} furnishes:
\begin{equation}
    \mathcal{L}\bm{s} = \Ke \dd \bm{e} + \Ken \dd \bm{\eta}
\label{eq:Lins}
\end{equation}
where:
\begin{subequations}
\begin{align}
\Ke &\coloneq \int \TotalStrainMatrix^{\top}\mathbb{C}\TotalStrainMatrix \dA  + \Kwg \label{eq:Kee} \\
\Ken &\coloneq \int \TotalStrainMatrix^{\top}\mathbb{C}\EnhanShapes\dA \label{eq:Ken} \\
\end{align}
\end{subequations}
and \(\Kwg\) is obtained by linearizing the Wagner term \eqref{eq:wagner-gl-strain-pi}:
\begin{equation}
\begin{aligned}
\Kwg =
\Twist & \SectionIntegral{
\left[\begin{matrix}
\mathbf{0} & r^{2}  \mathbb{C} \FrameBasis \otimes \FrameBasis & \mathbf{0} & \mathbf{0} \\
r^{2}  \FrameBasis \otimes \FrameBasis \mathbb{C} 
& r^{2}  \left(\bm{r}^{\times} \mathbb{C} \FrameBasis \otimes \FrameBasis - \FrameBasis \otimes \FrameBasis \mathbb{C} \bm{r}^{\times}\right) & r^{2}  \FrameBasis \otimes \FrameBasis \mathbb{C} \FrameBasis\otimes \WarpAxialShapes & r^{2}  \FrameBasis \otimes \FrameBasis \mathbb{C} \FrameBasis_{\beta} \otimes \nabla_{\beta} \WarpAxialShapes\\
\mathbf{0} & r^{2}  \WarpAxialShapes \otimes \FrameBasis \mathbb{C} \FrameBasis \otimes \FrameBasis & \mathbf{0} & \mathbf{0} \\
\mathbf{0} & r^{2}  \nabla_{\beta} \WarpAxialShapes \otimes \FrameBasis_{\beta} \mathbb{C} \FrameBasis \otimes \FrameBasis & \mathbf{0} & \mathbf{0} 
\end{matrix}\right]} \\
+\Twist^2 & \SectionIntegral{\left[\begin{matrix}
\mathbf{0} & \mathbf{0} & \mathbf{0} & \mathbf{0} \\
\mathbf{0} & r^{4} \FrameBasis \otimes \FrameBasis \mathbb{C} \FrameBasis \otimes \FrameBasis & \mathbf{0} & \mathbf{0} \\
\mathbf{0} & \mathbf{0} & \mathbf{0} & \mathbf{0} \\
\mathbf{0} & \mathbf{0} & \mathbf{0} & \mathbf{0}
\end{matrix}\right]
} \\
&+\SectionIntegral{\left[\begin{matrix}
\mathbf{0} & \mathbf{0} & \mathbf{0} & \mathbf{0} \\
\mathbf{0} & r^2 \AxialStress & \mathbf{0} & \mathbf{0} \\
\mathbf{0} & \mathbf{0} & \mathbf{0} & \mathbf{0} \\
\mathbf{0} & \mathbf{0} & \mathbf{0} & \mathbf{0} 
\end{matrix}\right]
}
\end{aligned}
\label{eq:section-wagner-tangent}
\end{equation}

Linearizing \eqref{eq:def-enhan-residual} furnishes:
\begin{equation}
\mathcal{L}\residn = \mathbf{K}_{e\eta}^{\top}\dd \bm{e} + \mathbf{K}_{\eta \eta}\dd \bm{\eta}
\qquad\implies\qquad 
\dd\bm{\eta} = \widetilde{\mathbf{K}} \dd \bm{e}
\label{eq:Linn}
\end{equation}
with
\begin{subequations}
\begin{align}
    \Knn &\coloneq \int \EnhanShapes^{\top}\mathbb{C}\EnhanShapes\dA \label{eq:Knn} \\
    \Kn &\coloneq -\Knn^{-1}\Ken^{\top} \label{eq:Kn}
\end{align}
\end{subequations}

From \eqref{eq:Lins} and \eqref{eq:Linn}:
\begin{SaveEquation}{eq:section-tangent}
\Ks = \Ke - \mathbf{K}_{\eta e}^{\top} \Kn
\end{SaveEquation}
